\subsection{Currencies in global macro}

\begin{qfdefinition}[The dollar in global macro: FX as shock absorber, transmission channel and expression; the dollar's share of turnover and reserves; the dollar index and its weights; and the SDR]\label{def:us-dollar-dxy-sdr}
\qffield{Statement} Gliner: ``The U.S.\ dollar accounts for 85 percent of all foreign exchange (FX) transactions'', the shares summing to 200 percent because each trade involves two currencies (Figure~7.1, BIS Triennial Survey, April 2013), and more than 60 percent of central banks' reserves, declining gradually \citep[chap.~7, sec.\ ``The Role of the U.S.\ Dollar'', PDF pp.~180--181]{s7-gliner}. The U.S.\ Dollar Index is the dollar against a basket changed once, in 1999, with weights EUR 57.6, JPY 13.6, GBP 11.9, CAD 9.1, SEK 4.2 and CHF 3.6 percent, traded as the DX future and the UUP ETF (p.~181).

The SDR was created by the IMF in 1969, initially defined as 0.888671 grams of fine gold, later as a basket of currencies; it ``is neither a currency, nor a claim on the IMF'' and serves as the IMF's unit of account (sec.\ ``Special Drawing Rights (SDRs)'', pp.~187--188). The IMF's current factsheet gives the basket as the U.S.\ dollar, euro, Chinese renminbi, Japanese yen and pound sterling, with the valuation reviewed every five years and the current weights effective 1 August 2022 \citep[factsheet dated 7 October 2023]{imf-sdr-factsheet}.

\qffield{Notation} DXY the ICE U.S.\ Dollar Index, a geometric average of six exchange rates; SDR Special Drawing Right.

\qffield{Conditions} Gliner's 85 percent means the dollar is on one side of 85 percent of trades; he sources it to the BIS 2013 survey (Figure~7.1). In April 2025 the dollar ``was on one side of 89.2\% of all trades'', up from 88.4 percent in 2022, with the euro at 28.9, the yen at 16.8, sterling at 10.2 and the renminbi at 8.5 percent, out of a total of 200 \citep[preliminary results, sec.~2]{bis-triennial-2025}.

\emph{Index methodology}. ICE describes the index as ``a geometrically-averaged calculation of six currencies'', computed as $\text{USDX}=50.14348112\times\text{EURUSD}^{-0.576}\times\text{USDJPY}^{0.136}\times\text{GBPUSD}^{-0.119}\times\text{USDCAD}^{0.091}\times\text{USDSEK}^{0.042}\times\text{USDCHF}^{0.036}$, the exponent negative where the dollar is the quote currency; it ``originally was developed by the U.S.\ Federal Reserve in 1973'' as a bilateral trade-weighted average, futures have traded since 1985, and the weights match Gliner's \citep[FAQ, June 2015, p.~2]{ice-dxy-faq}; later changes, if any, are not reflected. The weights are fixed, changed only once, in 1999 (Gliner); the index is therefore not a trade-weighted index against today's partners: it gives the euro over half the weight and omits the renminbi and the Mexican peso.

\qffield{Practice} FX plays three roles in macro: shock absorber, transmission channel and expression instrument. That global macro is de facto USD-centric, that many local trades are disguised USD liquidity trades, and that stress appears first in the dollar index, the cross-currency basis and EM currencies are interpretations, not findings of Gliner's chapter.

\qffield{Pitfalls} The SDR is not a reserve currency anyone trades; official reserve diversification is tracked by the IMF's COFER data on reserve composition (Gliner's Figure~7.2), not from the SDR basket. Weights in percent (57.6 for the euro) are exponents of a geometric average, not shares of a basket of currency amounts; a 1 percent rise in EUR/USD lowers the index by about 0.576 percent.

\qffield{See also} \hyperref[def:fx-market-structure]{spot FX market structure}; \hyperref[def:currency-regimes-intervention]{currency regimes}; \hyperref[mod:fx-carry-trade]{FX carry}; \hyperref[meth:fx-valuation-signals]{FX valuation signals}.

\qffield{Further reading} The dollar's role and safe-asset demand \citep{wide-jkl-safe-asset}.

\end{qfdefinition}

\begin{qfmethod}[Gliner's FX valuation signals: equity performance, sovereign credit, risk reversals, trade and the trade-weighted index, export-partner growth, debt and the current account, reserves, PPP and GDP per capita, and policy-rate and LIBOR differentials]\label{meth:fx-valuation-signals}
\qffield{Statement} Gliner lists factors for a currency model \citep[chap.~7, sec.\ ``Valuation Techniques for Foreign Exchange'', PDF pp.~195--209]{s7-gliner}: relative \emph{equity index performance}, through flows (pp.~195--196); \emph{credit}, five-year sovereign CDS (Australia against Mexico, p.~197); \emph{sentiment}, ``call-put skew on three-month risk reversals'', to be ``back-tested on a currency-by-currency basis'' and correlating better in distress than in calm periods (p.~197); the \emph{trade} balance and current account (pp.~198--199); the \emph{trade-weighted index} (pp.~199--200); and \emph{economic activity} (p.~200).

His further factors are \emph{export-partner growth} weighted by export shares, with Australia's exports to China, Japan, Korea, India and Taiwan (pp.~200--201); \emph{debt to GDP} (p.~202); \emph{absolute and relative PPP}, with the Big Mac index, and Meese and Rogoff's finding that it does not forecast in the short run (pp.~203--205); \emph{GDP per capita} (p.~205); and \emph{carry} indicators, the policy-rate (OIS) differential and the LIBOR differential, decomposed as LIBOR $=$ fed funds $+$ (LIBOR $-$ fed funds) (pp.~206--208), with risk-adjusted carry (\hyperref[mod:fx-carry-trade]{FX carry}).

\qffield{Notation} Risk reversal: implied volatility of an out-of-the-money call minus that of an equally out-of-the-money put (usually 25-delta); TWI: a trade-weighted index of the currency against its partners.

\qffield{Conditions} Gliner's examples as of 2012--2013; forecasts to 2021--2025 by currency, outcomes to 2026. The list is a set of candidate factors with charts, not a tested model: Gliner gives no out-of-sample results, and whether such factors beat a random walk is the open question. Several cut both ways: strong activity can attract inflows and also worsen the trade balance.

\qffield{Practice} Build each factor from data available at the time and test against a random walk with economic as well as statistical criteria. The factors can be grouped into risk appetite, trade, sustainability, long-term anchors and carry; FX reserve adequacy (months of import cover, short-term debt coverage) fits beside them, although Gliner treats it under reserves elsewhere.

\qffield{Pitfalls} Gliner's account of the \emph{Balassa--Samuelson effect}, that as a country becomes richer ``foreign investment flows into the country'' and strengthens the currency (p.~205), is not the effect: Balassa--Samuelson says faster productivity growth in traded goods raises the relative price of non-traded goods and so the real exchange rate. GDP per capita adjusted for PPP corrects for price levels, not for inflation as he puts it. Meese and Rogoff's 1983 result concerned structural exchange-rate models against a random walk at short horizons, a wider statement than ``PPP is not an accurate forecasting tool''. LIBOR differentials are historical (\hyperref[def:money-market-reference-rates]{reference rates}). \emph{Own out-of-sample test of one factor, relative PPP}. The hypothesis and pass rules were fixed in the script before its first run but were not preregistered.
\begin{itemize}
\item \emph{Design.} Nine currencies against the dollar (JPY, GBP, CHF, CAD, AUD, NZD, SEK, NOK, DKK), on FRED's monthly H.10 \emph{averages}. The signal $z_t$ is the deviation of the log real rate from its expanding mean.
\item \emph{Timing.} A month is eligible as a predictor only if its price levels were known at $t$: monthly CPI exactly two months old, or quarterly CPI four to six months past its quarter end, never older. Outcomes need only spot rates, not future CPI.
\item \emph{Forecasts.} Expanding regressions of the next 1- and 12-month spot change on $z_t$, using only outcomes already observed, were compared with a random walk from January 1985. The effective sample ends are August 2021 (JPY), May--June 2025 (the monthly-CPI currencies) and September 2025 (AUD, NZD).
\item \emph{Results.} Full-sample slopes are negative for all nine at 12 months, as PPP predicts. The forecasts lose to the random walk out of sample for every currency: pooled $R^2=-2.0$ percent at 1 month and $-14.9$ percent at 12 (\hyperref[meth:nested-forecast-tests]{Clark--West} $t=-1.3$ and $-0.2$).
\item \emph{Portfolio.} A monthly portfolio long the three cheapest and short the three dearest currencies returned $+0.09$ percent a year in spot terms, 1985 to May 2025 (Sharpe 0.01). Adding three-month interbank differentials gives $-0.17$ percent a year (2002--2025), before an assumed 5~bp per unit of turnover. Because it is computed from monthly averages, this is a statistical portfolio, not an executable trading result, and interbank rates are not the trader's funding or forward points.
\end{itemize}
Neither pass rule is met. \emph{A second factor: carry}. Gliner's carry indicator is the short-rate differential (pp.~206--208). It was tested on the same FRED files with a protocol written before any carry result was computed. The protocol is not blind to the data, which the PPP test had used. The predictor is the previous month's three-month interbank differential against the dollar. The model is a pooled expanding regression of 1- and 12-month dollar exchange-rate changes on it, forecasting from January 1995.
\begin{itemize}
\item \emph{Forecasts: not supported.} The pooled out-of-sample $R^2$ against a random walk is $-0.4$ percent at 1 month (Clark--West $t=-0.41$) and $-6.7$ percent at 12 (Clark--West $t=-1.97$), negative for every currency at 12 months. \emph{Panel inference}: Newey--West lags must be counted in calendar months, not in stacked currency rows, since same-month currencies are contemporaneously dependent. The Clark--West scores keep their forecast-origin month, are summed across currencies within it (each currency-month keeps weight $1/N$, so the pooled $R^2$ is unchanged), and the long-run variance is taken over 11 calendar months (12 as a sensitivity: $t=-1.94$). The forecasts still lose to the random walk.

The full-sample pooled slope is $+0.0017$ per percentage point at 12 months, against $-0.01$ under uncovered interest parity: high-rate currencies did not depreciate by the differential, on average. That is the forward-premium pattern, though it does not help forecast spot rates.
\item \emph{Portfolio: supported, narrowly.} Long the currencies above the cross-sectional median differential and short those below, earning spot change plus the interest accrual, net of 5~bp per unit of turnover. January 1985 to August 2026, the portfolio returned $+2.07$ percent a year (volatility 5.6 percent, Sharpe 0.37, Newey--West $t=2.04$, against the threshold of~2). Its three worst months were March 1985, September 2008 and February 2020. The portfolio has one return a month, so its Newey--West $t$ uses 12 monthly lags.
\end{itemize}
So carry pays through the accrued differential, not through forecastable exchange-rate moves. The portfolio is a statistical return on monthly-average prices and interbank rates, not an executable trade: an executable return would use month-end (or trade-time) rates, forward points or funding actually paid, and spreads; none of these was measured. Its $t$ barely clears the rule. Limits: monthly averages rather than month-end rates; revised CPI, not vintages; dollar bases only; interbank rates in place of policy-rate or OIS differentials; two of Gliner's twelve factors, since the others need data (CDS, risk reversals, export shares) not used here.

\qffield{See also} \hyperref[mod:fx-carry-trade]{FX carry}; \hyperref[def:macroeconomic-indicators]{external accounts}; \hyperref[def:sovereign-credit-defaults]{sovereign credit}.

\qffield{Further reading} Order flow and fundamentals in exchange-rate determination, and currency value \citep{wide-menkhoff-value}.

\end{qfmethod}
